\originalsection{位势流的拓扑与理想球体流}
无旋流局部有速度势, 全域单值势还需检查区域的拓扑与环量. 原课几个基本解显示这一差别. 二维点源$\phi=A\log r$给$v_r=A/r$, 单位深度流量$2\pi A$; 三维点源$\phi=-A/r$给$v_r=A/r^2$, 总流量$4\pi A$. 它们只在$r>0$满足无源Laplace方程. 二维势涡$\phi=\Gamma\theta/(2\pi)$给$v_\theta=\Gamma/(2\pi r)$, 在原点外旋度为零, 环量却为$\Gamma$. 势随绕原点一圈增加$\Gamma$, 因而不是全域单值函数.

对常密度无旋流, Euler方程给$\nabla(\phi_t+v^2/2+p/\rho+\Phi)=0$, 故括号等于只依赖$t$的函数$b(t)$. 可以给$\phi$加纯时间函数来改变$b$, 速度仍不变; 对空间有变化的压强却不能用这一自由度消去. 这就是非定常位势流的Bernoulli关系.

原课的三维例子是半径$a$固定球体, 远处流体速度$U\hat{\vect z}$. 方程为$\nabla^2\phi=0$, 球面只要求不可穿透$\partial_r\phi|_a=0$, 远处$\phi\sim Ur\cos\theta$. 由于边界条件具有$\cos\theta$的角依赖, 选轴对称偶极解$\phi=U(r+A/r^2)\cos\theta$. 求径向导数并代入球面, 得$A=a^3/2$, 因而
\[
v_r=U(1-a^3/r^3)\cos\theta,\qquad
v_\theta=-U(1+a^3/(2r^3))\sin\theta.
\]
球面切向速度为$-3U\sin\theta/2$, 伯努利式给
\[
p(a,\theta)=p_\infty+\frac{\rho U^2}{8}(9\cos^2\theta-5).
\]
压强在$\theta$与$\pi-\theta$相同, 沿流向的压力积分为零. 这是d'Alembert零阻力悖论的球体实例. 现实中的无滑动条件不在这个理想模型中, 增大Re也不意味着可在壁面处随意去掉黏性.

\originalsection{声速、Mach数与扰动传播}
前述声波中的$c_s^2$可写为体积模量$B$除以平衡密度: $B=\rho_0(\partial p/\partial\rho)_s$, $c_s^2=B/\rho_0$. 线性Euler式还给$\partial_t(\nabla\times\vect v)=0$; 如果扰动初始无旋, 其速度同样满足波动方程. 压强、密度和纵向速度波由连续性及状态方程联系, 不是三个任意独立波形.

特征流速$U$与声速之比$\mathrm{Ma}=U/c_s$称Mach数. 低Mach、弱密度变化且无强热源的条件下, 声传播比流动变化快, 常能采用不可压缩近似. 单凭$\mathrm{Ma}\ll1$并不排除由温度或分层导致的密度差. 对匀速超声源, 时刻$t-\tau$发出的球形扰动半径为$c_s\tau$, 源走过距离为$U\tau$; 包络锥半角$\gamma$满足$\sin\gamma=c_s/U=1/\mathrm{Ma}$. 原课用这一几何图说明超声扰动不能传到上游. Mach锥的线性传播几何本身不是有限振幅冲击波跃变条件的证明.
